\section{Discussion: Putting the Mapping to Work}
\label{sec:discussion}

The value of the methods\,$\times$\,requirements matrix lies less in its rows and columns than in how a practitioner uses them to convert scattered evaluation activity into structured conformity evidence. This section translates the mapping into an actionable playbook, illustrates it on a worked example, situates it within a US-led assurance ecosystem, and explains why a single evidence base can serve all three frameworks.

\subsection{A Conformity-Evidence Playbook}
We propose a four-step routine for assembling an evidence dossier directly from the matrix. \textit{First}, scope the system: classify it against the obligation triggers of each framework, namely the high-risk categories of the EU AI Act \cite{euaiact2024}, the outcome categories of the NIST AI RMF \cite{nist2023airmf}, and the management-system controls of ISO/IEC 42001 \cite{iso42001}, to determine which columns of Table~\ref{tab:matrix} are in scope. \textit{Second}, read down each in-scope column to enumerate the technical evaluation methods that can supply evidence for that requirement; the matrix thus acts as a requirements-to-tests index rather than a literature taxonomy. \textit{Third}, execute the selected methods and capture their outputs as durable artifacts (model cards \cite{puhlfurss2025modelcards}, system cards \cite{reddy2025systemcards}, audit cards \cite{staufer2025auditcards}, and technical documentation templates \cite{lucaj2025techops}), so that each evaluation produces a citable record rather than a transient result. \textit{Fourth}, reconcile coverage using the checklist in Table~\ref{tab:checklist}, which flags requirements with no associated evaluation evidence and exposes residual gaps before an assessor does. The dossier that emerges is a traceable chain from obligation to method to artifact, which is precisely the structure that certification proposals \cite{pan2024modelreporting} and automated compliance frameworks \cite{compliance2025pasta} assume but rarely make explicit.

\subsection{Worked Mini-Example}
Consider a hiring-screening model, which the EU AI Act treats as high-risk, triggering data-governance and accuracy-and-robustness duties \cite{euaiact2024}. Following the playbook, the in-scope columns include EU Art.~10 (data and data governance) and Art.~15 (accuracy, robustness, and cybersecurity). Reading down these columns, the practitioner selects dataset documentation and representativeness analysis to evidence Art.~10, and disaggregated accuracy testing, subgroup fairness evaluation, and adversarial robustness testing to evidence Art.~15; the fairness evaluations also speak to the non-discrimination questions that recent work shows are only partially captured by the Act's accuracy provisions \cite{meding2025algorithmicdiscrimination}. The same evaluations populate the NIST \textsc{Measure} function, whose analysis, benchmarking, and monitoring outcomes are evidenced by the identical test artifacts \cite{nist2023airmf}, and they supply evidence toward ISO/IEC 42001 controls A.6 (AI system life-cycle) and A.7 (data for AI systems) by producing records of documented, repeatable processes \cite{iso42001}. A clinical-decision model follows the same path with domain-specific robustness and monitoring evidence, illustrating how the matrix is intended to apply across high-risk application domains rather than encoding one sector's assumptions; we offer these as worked illustrations, not as a validation study.

\subsection{From a Test List to an Assurance Case}
A list of relevant tests is not yet an assurance case. To show how the crosswalk feeds one, we sketch a single assurance argument for the hiring model's Article~15 accuracy obligation, with the elements an assessor expects. \textit{Obligation:} EU AI Act Article~15 (appropriate accuracy). \textit{Scoped claim:} ``the system's selection accuracy and error rates are adequate and stable for the stated intended purpose and population.'' \textit{System context:} defined intended purpose, candidate population, and operating constraints. \textit{Evaluation protocol:} disaggregated accuracy and error-rate testing on a held-out, representative sample, with distribution-shift checks. \textit{Metric and acceptance criterion:} false-negative-rate parity within a pre-registered tolerance across protected groups, and overall accuracy above a domain-justified threshold; the criterion, not merely the test, is what discharges the claim. \textit{Result and uncertainty:} measured rates with confidence intervals. \textit{Residual risk:} performance on under-represented subgroups and post-deployment drift. \textit{Owner and cadence:} a named accountable role, with re-evaluation on each model or data change. The crosswalk locates \emph{which} evaluations are relevant; the assurance case adds the claim, context, threshold, uncertainty, residual risk, and sign-off that convert a relevant test into admissible evidence. This is the template practitioners should instantiate per obligation, and the layer a purely method-centric reading of the matrix omits.

\subsection{Evidence Levels and Stakeholders}
Two cautions keep the playbook from collapsing into checklist compliance. First, the \emph{unit of evaluation} matters: the same family yields evidence at the model level (e.g., a benchmark score), the system level (the model within its guardrails and tools), the workflow level (the human--AI process), and the organization level (governance and monitoring). High-risk conformity depends on the higher levels, yet most cited methods operate at the model level; practitioners should record, for each artifact, which level it actually evidences, and treat a model-level result as insufficient for a system- or workflow-level obligation. Second, conformity evidence implicates a \emph{stakeholder set} wider than the provider: deployers and operators (who run human oversight), affected persons (whose recourse and complaints can serve as evidence of oversight and fairness in practice), auditors and notified bodies (who judge sufficiency), regulators, and downstream integrators and third-party model vendors (whose dependencies must be represented). Human-oversight evidence in particular should distinguish the mere availability of a human from measured \emph{effectiveness} (detection, appropriate reliance, intervention success, escalation latency, authority, workload, training, and contestability), and should, where the obligation concerns affected persons, include evidence from operators and affected parties rather than from the model alone.

\subsection{US National-Importance Framing}
For a US audience, the playbook is most naturally anchored on the NIST AI RMF as the organizing spine, with the \textsc{Measure} function and the Generative AI Profile \cite{nist2024genai} supplying the evaluation vocabulary into which EU and ISO obligations are then translated. This NIST-centric framing matters for national competitiveness: a voluntary, outcome-based core lets US developers build one evaluation pipeline that also supports conformity evidence for export to regulated markets, an asymmetry that cross-regional studies highlight between voluntary and prescriptive regimes \cite{almaamari2025crossregional}. The framing also creates the demand-side conditions for an independent assurance market, in which third-party auditors verify evaluation claims rather than merely accept self-attestation \cite{auditing2024frontier}, and in which continuous-assurance postures extend point-in-time conformity to systems that evolve after deployment \cite{lupo2026trustworthyposture}.

\subsection{One Evidence Base, Three Frameworks}
The deeper implication of Table~\ref{tab:matrix} is that the three frameworks are largely non-redundant \textit{in obligation} yet highly overlapping \textit{in evidence}: in our mapping, all twelve method families have primary or supporting relevance to requirements in all three frameworks (Table~\ref{tab:matrix}), so no family is specific to a single regime. Because a single robustness test or fairness audit can be cited against a NIST outcome, an EU article, and an ISO control simultaneously, the marginal cost of multi-framework conformity is far lower than the number of frameworks would suggest, an economy that crosswalk efforts between the AI Act and international standards have begun to formalize \cite{hernandez2024knowledgegraph}, and that ISO/IEC 23894 risk guidance reinforces \cite{iso23894}. The practitioner's task therefore shifts from running parallel compliance programs to maintaining one well-documented evaluation base and re-indexing its artifacts against each framework's vocabulary, which is the central operational claim this review advances.
